\subsection{Whole-archive degradation and task-level uncertainty}
The whole-archive benchmark jointly masks selected gaps in several years of a target station before fitting its daily models, climatology or residual histories. Unmasked target observations and neighboring physical stations remain available in this experiment. The benchmark contains 42 stations with at least twenty complete reference years in seven regions. One Finnish record ends in 1994, leaving 41 stations for forward evaluation. Twenty paired mask repetitions use both empirical calendars and controlled 30--90-day gaps; after seven observed burn-in years, each reference year is degraded with probability 0.35. These 3,320 archive experiments do not reuse any masked target day when fitting another year. The complete reference, including those hidden days, is used for scoring; its unmasked values can also enter local fitting. The archived single-year plug-in is retained as an optimistic-information paired comparator on identical masks.

Offline daily fits use the remaining visible history. Forward annual models and interval offsets use source observations before 2010 and exclude the evaluated test stations and their connected groups. The original partitioned role assignment can retain other training-only groups in the same named region: this occurs in the HUC17 estimator and in NVE interval calibration. Accordingly this forward experiment is held-test-group, not uniformly held-region, and is applied to 2010 onward. Forward local daily fitting and dependence estimation also exclude water observations from 2010 onward. These are current-vintage retrospective forward experiments, not a historical real-time field trial. Annual regression weights remain source-only and use source-fit seed zero throughout. Source observations do not change across archive masks, so those fits are reused; every local stage that uses the degraded target archive is refitted. Twenty archive masks are not twenty new global model fits.

On each degraded archive, the retention set (90\% warm-season rule or all supported years) is crossed with the annual estimator (daily plug-in, equally calibrated stochastic regression or D2). Every slope is scored against the same complete-record slope. For the four MSE cells $A$ (rule/plug-in), $B$ (all/plug-in), $C$ (rule/alternative) and $D$ (all/alternative), the symmetric retention contribution is $\{(A-B)+(C-D)\}/2$, and the estimator contribution is $\{(A-C)+(B-D)\}/2$. They sum to $A-D$; the interaction is $(A-B)-(C-D)$. Complete paired cells, rather than different available subsets, determine each decomposition.

We refit both methods' annual quantile offsets on their original disjoint source calibration groups after imposing the observed-window lower bound. Both distributions and their samples respect the bound when one is available. Seven stored quantiles identify only the central 90\% of the distribution. We extrapolate both outer tails using the adjacent normal-quantile spacing and project samples onto the bound; the endpoint-constant alternative is retained on paired cases. The normal extrapolation is an assumption, not a validated tail law. We use 256 stratified samples per year and compare independent ranks, auxiliary shared year-bootstrap regression states, and three-year residual blocks on exactly the same annual marginals. Shared states come from a small harmonic/air Ridge model, not the D2 tree posterior or the selected donor model. These couplings are dependence sensitivities, not posterior samples or coverage theorems.

The task-level comparator calibrates absolute slope error divided by an independent-draw interval half-width. Every source connected group receives equal total weight in the empirical 95th percentile. Calibration uses only unused rule-calibration groups, re-predicted with the held-region model; forward errors end before 2010. We withhold the calibration when fewer than three nondegenerate source groups remain, the calibration factor is effectively zero, or the target half-width lies outside the source 1st--99th percentile range. These thresholds are empirical support checks rather than a finite-sample guarantee. Source support, availability and matched-case width, coverage and interval score are reported. All uncertainty scores target the finite complete-record slope, not climate-trend uncertainty.
